% SPM-Kit Fathom Usage Guide --- source 0.1.5.dev0 | GitHub 0.1.4 | PyPI 0.1.2
\documentclass[a4paper,11pt]{article}

% == Packages ==
\usepackage[utf8]{inputenc}
\usepackage[T1]{fontenc}
\usepackage{amsmath}
\usepackage[margin=2.2cm, bottom=2.5cm]{geometry}
\usepackage{microtype}
\usepackage[colorlinks=true, linkcolor=blue!60!black, urlcolor=blue!60!black, citecolor=blue!60!black]{hyperref}
\usepackage[table]{xcolor}
\usepackage{booktabs}
\usepackage{tabularx}
\usepackage{longtable}
\usepackage{array}
\usepackage{enumitem}
\usepackage[breakable, skins]{tcolorbox}
\usepackage{fancyhdr}
\usepackage{lastpage}
\usepackage{upquote}
\usepackage{textcomp}
\usepackage{xurl}
\usepackage{calc}
\usepackage{needspace}
\usepackage{graphicx}
\usepackage{listings}

% == Colours ==
\definecolor{noteblue}{RGB}{30,100,180}
\definecolor{warnamber}{RGB}{180,120,20}
\definecolor{tipgreen}{RGB}{30,130,70}
\definecolor{codegray}{RGB}{50,50,50}
\definecolor{lightgray}{RGB}{245,245,245}

% == Code block style ==
\newcommand{\cmd}[1]{\texttt{\small #1}}

\lstset{
  basicstyle=\small\ttfamily,
  backgroundcolor=\color{lightgray},
  frame=none,
  breaklines=true,
  breakatwhitespace=false,
  postbreak=\mbox{\textcolor{gray}{$\hookrightarrow$}\space},
  showstringspaces=false,
  aboveskip=4pt,
  belowskip=4pt,
  xleftmargin=4pt,
  xrightmargin=4pt,
  framexleftmargin=4pt,
  framexrightmargin=4pt,
  columns=flexible,
  keepspaces=true,
  upquote=true,
}



% == Note / warning boxes ==
\newtcolorbox{infobox}[1][]{colback=noteblue!5, colframe=noteblue!40,
  title={\textcolor{noteblue}{#1}}, fonttitle=\bfseries, breakable}

\newtcolorbox{warnbox}[1][]{colback=warnamber!5, colframe=warnamber!50,
  title={\textcolor{warnamber}{#1}}, fonttitle=\bfseries, breakable}

% == Page style ==
\pagestyle{fancy}
\fancyhf{}
\fancyhead[L]{\small SPM-Kit {\textperiodcentered} Fathom --- source 0.1.5.dev0}
\fancyhead[R]{\small\thepage/\pageref{LastPage}}
\renewcommand{\headrulewidth}{0.4pt}

% == Captionless table continuation ==
\newcommand{\tcont}[1]{%
  \multicolumn{1}{c}{\small\dots\ \emph{#1 continued}}\\}

\begin{document}

% ===========================================================
% TITLE
% ===========================================================

\begin{titlepage}
\raggedright
\vspace*{4cm}

{\Huge\bfseries SPM-Kit {\textperiodcentered} Fathom}\\[0.4cm]
{\LARGE Usage Guide}\\[1.5cm]
{\large Source manual \texttt{0.1.5.dev0} \,(Alpha)}\\[0.3cm]
{\large 2026-07-29}\\[0.3cm]
{\large GitHub release \texttt{0.1.4} {\textperiodcentered} PyPI \texttt{0.1.2}}\\[2cm]

{\normalsize
  \textbf{Creator, author, and lead developer:} Jos\'e Labarca Baeza\\
  \textbf{License:} MIT\\
  \textbf{Repository:} \url{https://github.com/kegouro/spmkit}\\
  \textbf{Site:} \url{https://kegouro.github.io/spmkit}\\[1cm]

  \emph{Canonical documentation: This PDF is a printable companion. The Markdown
  version at \texttt{docs/user-guide.md} is the authoritative reference and may
  contain additional detail. Both are verified against the development tree
  \texttt{0.1.5.dev0}. The latest tagged source release is \texttt{0.1.4}; PyPI
  currently serves \texttt{0.1.2}.}
  \\[0.6cm]
  \emph{Acknowledgements:} Tom\'as Corrales and the SPM Lab at Universidad T\'ecnica
  Federico Santa Mar\'ia provided selected experimental datasets and laboratory context
  during the development and evaluation of SPM-Kit. Mar\'ia Saavedra Fredes and Benjamin
  Schleyer helped locate and share candidate datasets for the validation campaigns.
}

\vspace{1cm}
{\normalsize
  SPM-Kit is an open-source Python toolkit for analysing AFM/KPFM scanning-probe-microscopy
  data. Fathom is its desktop workspace. This guide covers installation, conceptual model,
  supported formats, the Fathom GUI, CLI reference, Python API, scientific validation,
  and extensibility.
}

\vfill
{\footnotesize
  This open-source project has been developed without dedicated institutional funding.
  Development status: Alpha (3~---~Alpha). API may change before 1.0.
}
\end{titlepage}

\tableofcontents
\clearpage

% ===========================================================
% 1. INTRODUCTION
% ===========================================================
\section{Introduction}

\subsection{What is SPM-Kit?}

SPM-Kit is an open-source Python toolkit for analysing AFM/KPFM scanning-probe-microscopy data.
It reads demonstrated instrument-file variants and produces inspectable results: roughness
(ISO~25178), KPFM contact-potential statistics, nanomechanical property maps,
thermal-tune cantilever calibration, grain detection, spectral/fractal analysis, and
single-molecule force spectroscopy.

SPM-Kit was created and is authored and led by Jos\'e Labarca Baeza as an independent
software project. Acknowledged dataset and laboratory-context contributions do not create
software authorship or institutional ownership.

\subsection{What is Fathom?}

Fathom is the desktop workspace application bundled with SPM-Kit. It is a PyQt6~+~pyqtgraph
GUI organised into \emph{perspectives} (task-oriented views), not flat tabs. A command
palette (\texttt{Ctrl+K}) provides keyboard access to every function.

\subsection{Component relationship}

\begin{lstlisting}[language=bash]
spmkit (Python package)
  spmkit.core   -- numerical engine (pure Python, no UI imports)
  spmkit.cli    -- command-line interface (Typer + Rich)
  spmkit.gui    -- Fathom desktop workspace (PyQt6 + pyqtgraph)
\end{lstlisting}

\begin{itemize}
  \item \texttt{spmkit} is the Python package and CLI entry point.
  \item \textbf{Fathom} is the graphical workspace.
  \item \textbf{SPM-Kit} is the project/ecosystem name.
\end{itemize}

\subsection{What SPM-Kit does \emph{not} do}
\begin{itemize}[nosep]
  \item It is \textbf{not a microscope controller}.
  \item It is \textbf{not a certified metrology tool}.
  \item It does \textbf{not} provide medical, clinical, or regulatory reports.
  \item Simulation features are \textbf{educational}, not quantitative digital twins.
\end{itemize}

\subsection{Status}
SPM-Kit is \textbf{alpha-stage} (Development Status~3). The strongest retained evidence is
\textbf{LEVEL~3} for Sa, Sq and Sz under frozen campaign conditions: 48 synthetic matrices
produced 144/144 conforming comparisons and 12 public experimental GWY records produced
36/36 shared-matrix comparisons against Gwyddion~2.71. Six Nanoscope~III \texttt{.spm}
files support a partial \textbf{LEVEL~2} claim with a documented pre-freeze unblinding
incident. KPFM, spectral, grain, and resonance utilities remain Level~1; contact and chain
models have claim-scoped synthetic Level~2 evidence. No general Level~4 or Level~5 claim is
made (see \S\ref{sec:validation}).

% ===========================================================
% 2. INSTALLATION
% ===========================================================
\section{Installation}

Requirements: \textbf{Python\,{\textgreater}=\,3.11}.

The documented source is \texttt{0.1.5.dev0}; the latest GitHub release is \texttt{0.1.4};
PyPI currently serves \texttt{0.1.2}. Choose the source explicitly.

\subsection{Core only}
\begin{lstlisting}[language=bash]
python -m pip install "spmkit @ git+https://github.com/kegouro/spmkit@v0.1.4"
\end{lstlisting}

Installs the numerical core and CLI. Optional format dependencies still require extras.

\subsection{GUI / Fathom}
\begin{lstlisting}[language=bash]
python -m pip install "spmkit[gui] @ git+https://github.com/kegouro/spmkit@v0.1.4"
\end{lstlisting}

Adds: PyQt6, pyqtgraph, matplotlib, Crameri colormaps, matplotlib-scalebar.

\subsection{All features}
\begin{lstlisting}[language=bash]
python -m pip install "spmkit[all] @ git+https://github.com/kegouro/spmkit@v0.1.4"
\end{lstlisting}

\subsection{Extras matrix}

\begin{longtable}{@{}>{\ttfamily}l p{7cm} l@{}}
\toprule
\textnormal{Extra} & \textnormal{Provides} & \textnormal{Deps} \\
\midrule
gui    & Fathom desktop workspace & PyQt6, pyqtgraph, matplotlib \\
viz    & Publication figures, colormaps & matplotlib, cmcrameri, scalebar \\
gwy    & Gwyddion .gwy read/write & gwyfile \\
hdf5   & HDF5 import/export & h5py \\
grains & Grain/particle detection & scipy \\
report & HTML/PDF reports & Jinja2 + spmkit[viz] \\
nanosurf & Optional NSFopen dependency & NSFopen \\
afm    & JPK QI, .ibw, HDF5, NT-MDT & afmformats \\
jpk    & JPK TIFF force curves & tifffile \\
parallel & Parallel force-volume & joblib \\
pandas & DataFrame export & pandas \\
dev    & Lint, test, type-check & pytest, ruff, mypy, black \\
test-gui & GUI test runner & pytest-qt \\
docs   & Documentation build & mkdocs-material \\
all    & Reader, GUI, visualization, report, grain, HDF5 and GWY extras & --- \\
\bottomrule
\end{longtable}

\subsection{Development installation}
\begin{lstlisting}[language=bash]
git clone https://github.com/kegouro/spmkit
cd spmkit
pip install -e ".[all,dev]"
\end{lstlisting}

\subsection{Headless / HPC}
Install the tagged core source shown above. Add \texttt{[viz]} for headless figure
generation.

\subsection{Platform notes}
Developed and tested on \textbf{macOS} and \textbf{Linux}. Windows is not regularly tested.
Fathom requires a display; set \texttt{QT\_QPA\_PLATFORM=offscreen} for headless tests.

% ===========================================================
% 3. VERIFICATION
% ===========================================================
\section{Verifying the installation}

\begin{lstlisting}[language=bash]
spmkit --version          # 0.1.4 tag; 0.1.2 PyPI; 0.1.5.dev0 main
python -c "import spmkit; print(spmkit.__version__)"
\end{lstlisting}

% ===========================================================
% 4. CONCEPTUAL MODEL
% ===========================================================
\section{Conceptual model}

\subsection{Architecture}

Three strictly separated layers (enforced by \texttt{tests/test\_architecture.py}):

\begin{itemize}[nosep]
  \item \texttt{core/} --- pure Python, zero UI imports.
  \item \texttt{cli/} --- Typer + Rich. Only imports the public \texttt{core} API.
  \item \texttt{gui/} --- PyQt6 + pyqtgraph. Only imports the public \texttt{core} API.
\end{itemize}

\subsection{Data model}

\begin{lstlisting}[language=bash]
from spmkit import load, SPMData, SPMChannel

data: SPMData = load("scan.nid")
ch: SPMChannel = data["Z-Axis"]

ch.data         # 2D ndarray in physical units
ch.unit         # "m", "V", "{\textdegree}"
ch.x_range      # metres
ch.y_range      # metres
ch.direction    # "forward" | "backward"
\end{lstlisting}

\subsection{Force-spectroscopy domain}

For force curves: \texttt{ForceCurve}, \texttt{ForceSegment}, \texttt{ForceVolume}
(lazy-loaded), \texttt{Calibration}.

\subsection{Result objects}
Analysis results are immutable dataclasses. Use \texttt{to\_dict()} before passing
them to the export helpers.

% ===========================================================
% 5. SUPPORTED FILE FORMATS
% ===========================================================
\section{Supported file formats}

\begin{longtable}{@{}>{\raggedright\arraybackslash}p{2.8cm} p{3.0cm} p{2.5cm} p{2.0cm} p{3.8cm}@{}}
\toprule
Extension & Source & Type & Access & Status \\
\midrule
.nid & NanoSurf classic & Image + spectroscopy & Read & Implemented, scoped \\
.nhf & NanoSurf HDF5 & Images & Read & Experimental; hdf5 extra \\
.gwy & Gwyddion & Images & R/W & Implemented; gwy extra \\
.spm/.00N & Bruker / DI & Images & Read & \textbf{Partial LEVEL 2} \\
.jpk-force/.jpk & JPK & Single curve & Read & Implemented, scoped \\
tagged TIFF & JPK export & Force data & Read & Experimental; jpk extra \\
.jpk-qi-data/map/series & JPK & Force data & Read & afmformats adapter \\
.ibw/.h5/.tab & Adapter sources & Force data & Read & afmformats adapter \\
.npz & SPM-Kit Phantoms & Images & Read & Declared bundle only \\
\bottomrule
\end{longtable}

\subsection{Maturity vocabulary}
\begin{description}[nosep]
  \item[LEVEL 0] Claimed, without retained executable evidence.
  \item[LEVEL 1] Software verified under declared scope.
  \item[LEVEL 2] Known values recovered under declared numerical scope and tolerance.
  \item[LEVEL 3] Frozen external cross-validation under declared conditions.
  \item[LEVEL 4] Scoped physical-reference validation.
  \item[LEVEL 5] Independent reproduction of the declared result.
\end{description}

\subsection{What the .nid evidence means}
The parser has synthetic byte-budget and orientation tests plus selected laboratory-context
comparisons. The private instrument corpus is not distributed, so this supports scoped Level~1
software evidence rather than universal NanoSurf compatibility. The Sa/Sq/Sz shared-matrix
campaigns are separate Level~3 algorithm evidence. Parser, algorithm, and physical-validation
claims must remain distinct.

% ===========================================================
% 6. QUICK START
% ===========================================================
\section{Quick start}

\begin{lstlisting}[language=bash]
# Inspect a file
spmkit info scan.nid

# Roughness (ISO 25178)
spmkit roughness scan.nid -c Z-Axis --level plane

# Publication figure
spmkit figure scan.nid -c Z-Axis -o topo.png --colormap batlow

# Launch Fathom
spmkit gui scan.nid

# Force curve fit
spmkit forcecurve curva.jpk-force --model dmt --tip-radius 2e-8
\end{lstlisting}

\needspace{4\baselineskip}
Python:

\begin{lstlisting}[language=bash]
from spmkit import load
from spmkit.core.analysis import leveling, roughness

data = load("scan.nid")
flat = leveling.plane_fit(data["Z-Axis"])
stats = roughness.statistics(flat)
print(f"Sa = {stats.Sa:.3g} {stats.unit}")
\end{lstlisting}

% ===========================================================
% 7. FATHOM DESKTOP WORKSPACE
% ===========================================================
\section{Fathom desktop workspace}

\subsection{Launching}
\begin{lstlisting}[language=bash]
spmkit gui               # Fathom (default)
spmkit gui scan.nid      # open file on launch
spmkit gui --legacy      # legacy 7-tab app (fallback)
\end{lstlisting}

\subsection{Opening data}
Drag and drop a file onto the window, or \texttt{Ctrl+O}. Fathom inspects the file and
routes image channels to image perspectives and force/spectroscopy data to force perspectives.
If a file contains both, Fathom asks which to open.

\subsection{Projects (.spmproj)}
\texttt{Ctrl+S} saves the open file, pipeline parameters, and active perspective as a
human-editable YAML file.

\subsection{Appearance}
\begin{itemize}[nosep]
  \item \texttt{Ctrl+Shift+L} --- toggle light/dark theme.
  \item \texttt{Ctrl+Shift+A} --- appearance dialog (theme presets, accent colour, font size).
  \item Settings persist between sessions.
  \item The theme feeds the Qt app, pyqtgraph, and matplotlib simultaneously.
\end{itemize}

\subsection{Command palette}
\texttt{Ctrl+K}, then type a command name and press Enter. Indexes all perspectives and
registered commands.

% ===========================================================
% 8. PERSPECTIVES
% ===========================================================
\section{Perspectives}

Fathom presents 12 perspectives through 6 modules. Each perspective shows only the
panels relevant to its workflow.

\needspace{4\baselineskip}
\subsection{Image perspectives}

\begin{longtable}{@{}>{\ttfamily}l p{9cm}@{}}
\toprule
\multicolumn{2}{c}{\textbf{Module: image}}\\
\midrule
image & Channel viewer, leveling (plane/poly/align-rows), colormap, line profile ROI, roughness + KPFM analysis, profile export.\\
grains & Particle/grain detection on levelled topography; colour overlay, size statistics. Requires \texttt{grains} extra.\\
spectral & Radial PSD (log-log), fractal dimension~D, Hurst exponent~H, correlation length.\\
resonance & Thermal tune: SHO fit $\rightarrow$ f$_0$, Q, spring constant by equipartition.\\
evaporation & Mass sensing: load series of tuning .nid files $\rightarrow$ f(t), mass, d\textsuperscript{2} law fit.\\
\bottomrule
\end{longtable}

\needspace{4\baselineskip}
\subsection{Force perspectives}

\begin{longtable}{@{}>{\ttfamily}l p{9cm}@{}}
\toprule
\multicolumn{2}{c}{\textbf{Module: force}}\\
\midrule
force & Navigate curves, fit Hertz/DMT/Sneddon, inspect modulus $\pm$ 1$\sigma$, adhesion, dissipation, R\textsuperscript{2}. Editable pipeline, Monte Carlo uncertainty. Export scientific CSV.\\
smfs & Single-molecule force spectroscopy: prominence-based rupture detection, WLC/FJC polymer fits, contour-length histogram.\\
map & Force-volume $\rightarrow$ property maps (modulus, adhesion, contact, R\textsuperscript{2}) + histogram. Vectorised (CPU/GPU) or per-curve pipeline.\\
batch & Process folders of force curves $\rightarrow$ summary CSV.\\
\bottomrule
\end{longtable}

\needspace{4\baselineskip}
\subsection{Other perspectives}

\begin{longtable}{@{}>{\ttfamily}l p{9cm}@{}}
\toprule
\multicolumn{2}{c}{\textbf{Module: figure}}\\
\midrule
figure & WYSIWYG publication figure editor: title, axes, colormap, scale bar, draggable annotations. Export PNG/SVG/PDF.\\
\midrule
\multicolumn{2}{c}{\textbf{Module: view3d}}\\
\midrule
view3d & 3D topography surface with hillshade lighting, visual Z exaggeration.\\
\midrule
\multicolumn{2}{c}{\textbf{Module: simulator}}\\
\midrule
simulator & Educational cantilever digital twin: thermal noise spectrum, resonance shift with added mass.\\
\bottomrule
\end{longtable}

% ===========================================================
% 9. END-TO-END WORKFLOWS
% ===========================================================
\section{End-to-end workflows}

\subsection{Topography leveling and roughness}
\textbf{CLI:}
\begin{lstlisting}[language=bash]
spmkit roughness scan.nid -c Z-Axis --level plane
\end{lstlisting}

\textbf{Python:}
\begin{lstlisting}[language=bash]
from spmkit import load
from spmkit.core.analysis import leveling, roughness

data = load("scan.nid")
flat = leveling.plane_fit(data["Z-Axis"])
stats = roughness.statistics(flat)
print(stats.Sa, stats.Sq, stats.Sz, stats.unit)
\end{lstlisting}

\textbf{GUI:} Drag file $\rightarrow$ Image perspective $\rightarrow$ select channel $\rightarrow$ Plane Fit $\rightarrow$ read results from Analysis panel.

\begin{warnbox}[Caveats]
Roughness depends on scan size, pixel density, and leveling method. Always report these
alongside roughness values.
\end{warnbox}

\subsection{Force-curve contact mechanics}

\begin{lstlisting}[language=bash]
spmkit forcecurve sample.jpk-force --model dmt --tip-radius 2e-8
\end{lstlisting}

\textbf{Contact models:} \texttt{sphere} (Hertz), \texttt{paraboloid} (Sneddon paraboloid),
\texttt{cone} (Sneddon cone), \texttt{dmt} (Hertz + adhesion offset). The experimental JKR
model (\texttt{spmkit jkr}) fits adhesive contact but has not been validated against an
independent reference.

\begin{warnbox}[Caveats]
Modulus depends on tip radius (uncertain), contact model (approximation), and calibration.
Report all three alongside the result. JKR is experimental.
\end{warnbox}

\subsection{Force-volume property mapping}

\begin{lstlisting}[language=bash]
spmkit forcemap volume.nid --model dmt --fast -f maps.png
spmkit forceexport volume.nid -o ./results
\end{lstlisting}

Creates property maps (modulus, adhesion, contact point, R\textsuperscript{2}), histograms, and CSV exports.
Use \texttt{--pipeline} for variable-length curves (QI data),
\texttt{--parallel} for multi-core, \texttt{--backend gpu} for GPU (requires CuPy).

\subsection{Thermal tune}
Load a thermal spectrum in the Thermal Tune perspective or use:
\begin{lstlisting}[language=bash]
from spmkit.core.analysis.mechanics import thermal_spring_constant
\end{lstlisting}

\subsection{Evaporation / mass sensing}
\begin{lstlisting}[language=bash]
spmkit evaporation ./tuning/ -k 0.3 -o evap.csv
\end{lstlisting}

Tracks f(t) $\rightarrow$ mass $\rightarrow$ evaporation rate. Fits the d\textsuperscript{2} law for diffusion-limited evaporation.

\subsection{KPFM}
\begin{lstlisting}[language=bash]
spmkit analyze scan.nid --tip-wf 4.8 -o ./results
\end{lstlisting}

Reports mean CPD, standard deviation, and sample work function
($\Phi_{\text{sample}} = \Phi_{\text{tip}} - e\cdot\text{CPD}$).

\subsection{Other workflows}
\begin{lstlisting}[language=bash]
spmkit grains scan.nid --min-size 10           # grain detection
spmkit psd scan.nid -c Z-Axis                   # fractal analysis
spmkit batch ./measurements/ -o summary.csv     # batch processing
spmkit fbatch ./curves/ -o batch.csv            # force batch
spmkit figure scan.nid -o topo.png --colormap batlow  # figure export
\end{lstlisting}

% ===========================================================
% 10. COMMAND-LINE REFERENCE
% ===========================================================
\section{Command-line reference}

19 commands are registered in the Typer application.

\begin{longtable}{@{}>{\ttfamily}l p{8.5cm}@{}}
\toprule
Command & Description \\
\midrule
info & Show file metadata and channels \\
roughness & ISO 25178 parameters (Sa, Sq, Sz, Ssk, Sku) \\
psd & Spectral: fractal dimension, Hurst, correlation length \\
analyze & Full pipeline: roughness + KPFM $\rightarrow$ CSV + JSON \\
nanomech & Force-curve fit (Hertz/DMT/Sneddon) $\rightarrow$ Young's modulus \\
grains & Grain/particle detection and size statistics \\
batch & Process folder $\rightarrow$ summary CSV \\
evaporation & Mass sensing: f(t) $\rightarrow$ mass, evaporation rate, d\textsuperscript{2} law \\
figure & Publication figure: PNG/SVG/PDF with scale bar \\
convert & Format conversion (.nid $\rightarrow$ .gwy / .h5) \\
verify & Up to 8 structural/numeric .nid checks \\
gui & Launch Fathom (or legacy with --legacy) \\
workspace & Force-curve workspace (redesign) \\
forcecurve & Fit single force curve (JPK/NanoSurf) \\
forcemap & Force-volume $\rightarrow$ property maps \\
forcereport & Force-volume report (HTML/LaTeX/PDF) \\
forceexport & Full export bundle to folder \\
fbatch & Batch force processing $\rightarrow$ CSV \\
jkr & \textbf{[Experimental]} JKR adhesive contact fit \\
\bottomrule
\end{longtable}

\subsection{Common options}
\begin{itemize}[nosep]
  \item \texttt{-{}-channel / -c} --- channel name (default: \texttt{Z-Axis} for images,
    \texttt{Deflection} for force)
  \item \texttt{-{}-output / -o} --- output path
  \item \texttt{-{}-level / -l} --- leveling: \texttt{plane}, \texttt{poly}, \texttt{none}
  \item \texttt{-{}-model} --- contact model: \texttt{sphere}, \texttt{paraboloid}, \texttt{cone}, \texttt{dmt}
  \item \texttt{-{}-tip-radius} --- tip radius in metres (default: \texttt{1e-08})
\end{itemize}

Run \texttt{spmkit <command> --help} for full options.

% ===========================================================
% 11. PYTHON API
% ===========================================================
\section{Python API}

\subsection{Public imports}

\begin{lstlisting}[language=bash]
from spmkit import load, SPMData, SPMChannel, __version__
from spmkit.core.analysis import leveling, roughness, kpfm, mechanics
from spmkit.core.analysis import spectral, grains, resonance
from spmkit.core.export import to_csv, to_json, to_hdf5
from spmkit.core.viz import FigureSpec, save_figure
\end{lstlisting}

\subsection{Loading and inspecting}

\begin{lstlisting}[language=bash]
data = load("scan.nid")    # auto-detect format
print(data.names)          # channel names
ch = data["Z-Axis"]        # SPMChannel
\end{lstlisting}

\subsection{Analysis}

\begin{lstlisting}[language=bash]
flat = leveling.plane_fit(data["Z-Axis"])
stats = roughness.statistics(flat)
cpd = kpfm.statistics(data["CPD"], tip_work_function=4.8)
frac = spectral.fractal_dimension(flat)
grains_result = grains.detect(flat, min_size=4)
\end{lstlisting}

\subsection{Force spectroscopy}

\begin{lstlisting}[language=bash]
from spmkit.core.io.forceload import load_nid_force
from spmkit.core.analysis.mechanics import fit_hertz

volume = load_nid_force("sample.nid")
result = fit_hertz(volume.curves[0], tip_radius=10e-9, model="dmt")
print(result.young_modulus, result.r_squared)
\end{lstlisting}

\subsection{Export and figures}

\begin{lstlisting}[language=bash]
to_csv(stats, "roughness.csv")
to_json(stats, "roughness.json")
to_hdf5(data, "scan.h5")

spec = FigureSpec(title="AFM Topography", colormap="batlow")
save_figure(flat, spec, "topo.png")
\end{lstlisting}

\subsection{Recipes}

\begin{lstlisting}[language=bash]
from spmkit.core.pipeline import Recipe, Step, run

recipe = Recipe(steps=(
    Step(op="calibrate"),
    Step(op="find_contact_point"),
    Step(op="fit_elasticity", params={"model": "dmt", "tip_radius": 2e-8}),
))
_, ctx = run(recipe, curve)
print(ctx["young_modulus"], ctx["r_squared"])
\end{lstlisting}

% ===========================================================
% 12. KEYBOARD SHORTCUTS
% ===========================================================
\section{Keyboard shortcuts}

\begin{tabularx}{\linewidth}{@{}>{\ttfamily}l X@{}}
\toprule
Shortcut & Action \\
\midrule
Ctrl+K & Command palette \\
Ctrl+O & Open file \\
Ctrl+S & Save project \\
Ctrl+M & Calculate map \\
Ctrl+Shift+R & Generate report \\
Ctrl+Shift+L & Toggle light/dark theme \\
Ctrl+Shift+A & Customize appearance \\
Ctrl+Left / Ctrl+Right & Previous / next curve \\
Ctrl+1 ... Ctrl+5 & Tab switch (legacy app only) \\
\bottomrule
\end{tabularx}

% ===========================================================
% 13. EXPORTS AND PROVENANCE
% ===========================================================
\section{Exports and provenance}

\subsection{Export formats}

\begin{tabularx}{\linewidth}{@{}>{\ttfamily}l l X@{}}
\toprule
Format & Via & Notes \\
\midrule
CSV & \texttt{to\_csv()}, CLI & Units in headers; empty cells for NaN \\
JSON & \texttt{to\_json()}, CLI & Structured; programmatic consumption \\
HDF5 & \texttt{to\_hdf5()}, convert & Self-describing binary; requires hdf5 extra \\
.gwy & \texttt{spmkit convert} & Round-trip with Gwyddion \\
PNG/SVG/PDF & \texttt{spmkit figure} & Publication-quality figures \\
HTML/PDF & \texttt{spmkit forcereport} & Full force-volume report \\
YAML & \texttt{core.pipeline} & Reproducible pipeline recipe \\
\bottomrule
\end{tabularx}

\subsection{Provenance}
\begin{itemize}[nosep]
  \item \texttt{SPMData.source\_path} records the original file path.
  \item \texttt{SPMData.metadata} preserves instrument parameters.
  \item \texttt{spmkit verify} performs up to eight structural and numeric checks over
        the declared .nid header and binary layout; malformed files can stop early.
  \item YAML recipes capture the exact pipeline parameters.
  \item Exports include metadata when the format supports it.
\end{itemize}

% ===========================================================
% 14. SCIENTIFIC VALIDATION
% ===========================================================
\section{Scientific validation}
\label{sec:validation}

Evidence belongs to a claim, data family, version, preprocessing path, tolerance, and
campaign. A high level for one metric never transfers automatically to an adjacent reader,
model, or sample.

\subsection{Evidence-level vocabulary}
\begin{description}[nosep]
  \item[LEVEL 0 --- CLAIMED] Documented without retained executable evidence.
  \item[LEVEL 1 --- SOFTWARE VERIFIED] Automated tests exercise the declared behavior.
  \item[LEVEL 2 --- NUMERICALLY VERIFIED] Known values are recovered within a stated scope and tolerance.
  \item[LEVEL 3 --- CROSS VALIDATED] An external route is compared under a frozen protocol.
  \item[LEVEL 4 --- PHYSICALLY VALIDATED] A physical reference supports the scoped capability.
  \item[LEVEL 5 --- REPRODUCIBILITY VALIDATED] An independent party reproduced the result.
\end{description}

\subsection{Status summary}

{\small
\begin{longtable}{@{}p{4.0cm} p{2.6cm} p{2.2cm} p{5.1cm}@{}}
\toprule
Capability & Evidence & Status & Limitations \\
\midrule
Sa/Sq/Sz synthetic & 48 cases; 144/144 & \textbf{LEVEL 3} & Shared matrices; no preprocessing; three metrics \\
Sa/Sq/Sz public GWY & 12 records; 36/36 & \textbf{LEVEL 3} & Algorithm track; parser observations separate \\
Nanoscope III .spm & 6 files; 18/18 metrics & \textbf{LEVEL 2} & Partial; pre-freeze unblinding; no blind holdout \\
NanoSurf .nid mapping & Byte/orientation tests & \textbf{LEVEL 1} & Private corpus not distributed \\
Hertz/cone/DMT/maps & Synthetic recovery & \textbf{LEVEL 2} & No certified calibration or broad physical campaign \\
JKR/WLC/FJC/SLS & Synthetic recovery & \textbf{LEVEL 2} & JKR/SLS experimental; no physical campaign \\
KPFM, PSD, grains, resonance & Software tests & \textbf{LEVEL 1} & No frozen public physical-reference campaign \\
Fathom & Offscreen GUI tests & \textbf{LEVEL 1} & Qt/platform behavior varies \\
\bottomrule
\end{longtable}
}

\subsection{Gwyddion as reference}
Gwyddion is an external software route only where a campaign declares its version, matrix
handoff, wrapper, preprocessing, and tolerance. It is not universal ground truth. Retained
summaries live at \url{https://github.com/kegouro/spmkit-validation}. Passing them does not
constitute physical validation or feature parity.

\subsection{Simulator}
The cantilever simulator is \textbf{educational only}. It models a damped harmonic
oscillator with additive noise. Not suitable for calibration or metrology.

% ===========================================================
% 15. SAFETY AND INTERPRETATION
% ===========================================================
\section{Safety and interpretation warnings}

\begin{itemize}
  \item \textbf{Results require domain judgment.} High R\textsuperscript{2} does not guarantee correct model choice.
  \item \textbf{Calibration values matter.} Young's modulus depends on tip radius (uncertain),
    spring constant (calibration), and contact model (approximation).
  \item \textbf{Check units.} \texttt{SPMChannel.unit} tells the physical unit.
  \item \textbf{Experimental readers may misinterpret metadata.} Format parsing is
    best-effort for reverse-engineered formats.
  \item \textbf{Fitted parameters are not automatically physical truth.}
  \item \textbf{SPM-Kit is not a medical, clinical, regulatory, or safety-critical instrument.}
  \item \textbf{Preserve original data.} SPM-Kit reads without modifying files.
\end{itemize}

% ===========================================================
% 16. EXTENSIBILITY
% ===========================================================
\section{Extensibility}

\subsection{Entry-point groups}

\begin{tabularx}{\linewidth}{@{}>{\ttfamily}l X@{}}
\toprule
Group & Purpose \\
\midrule
spmkit.plugins.v1 & Readers, analyses, domains \\
spmkit.gui.modules & Fathom modules (perspectives + panels) \\
\bottomrule
\end{tabularx}

\subsection{Adding a format}
Implement the \texttt{Reader} Protocol (\texttt{inspect + load}) and register via
\texttt{spmkit.plugins.v1}.

\subsection{Adding a perspective}
Define a \texttt{ModuleSpec} with \texttt{PanelSpec} and \texttt{PerspectiveSpec},
register via \texttt{spmkit.gui.modules}. See \texttt{docs/extending.md} for full details.

% ===========================================================
% 17. TROUBLESHOOTING
% ===========================================================
\section{Troubleshooting}

\begin{description}[nosep, font=\normalfont\ttfamily]
  \item[Missing GUI extra] \hfill \\ \texttt{pip install "spmkit[gui]"}
  \item[Qt platform plugin error] \hfill \\ Install Qt dependencies or set \texttt{QT\_QPA\_PLATFORM=offscreen}.
  \item[Unsupported format] \hfill \\ Check extension matches \texttt{.nid/.nhf/.gwy/.jpk-force}. For experimental readers, install \texttt{afm} or \texttt{jpk}.
  \item[Missing optional dependency] \hfill \\ Install the relevant extra (e.g. \texttt{pip install "spmkit[grains]"}).
  \item[Large force-volume memory] \hfill \\ Use lazy loading or the \texttt{--fast} CPU vectorised path.
  \item[GPU fallback] \hfill \\ CuPy is not bundled. Install separately: \texttt{pip install cupy-cuda12x}.
  \item[Export failure] \hfill \\ Check write permissions. Use \texttt{--output} with a writable path.
\end{description}

% ===========================================================
% 18. DEVELOPMENT
% ===========================================================
\section{Development and quality}

\begin{lstlisting}[language=bash]
git clone https://github.com/kegouro/spmkit
cd spmkit
pip install -e ".[all,dev,test-gui]"
pre-commit install
\end{lstlisting}

\textbf{Quality checks (CI gate):}
\begin{lstlisting}[language=bash]
make check        # lint + types + tests
make lint         # ruff check
make type         # mypy (strict on core)
make test         # pytest with coverage
\end{lstlisting}

\textbf{GUI tests} (headless):
\begin{lstlisting}[language=bash]
QT_QPA_PLATFORM=offscreen pytest tests/gui -q
\end{lstlisting}

\textbf{Documentation:}
\begin{lstlisting}[language=bash]
pip install -e ".[docs]"
python -m mkdocs build --strict
\end{lstlisting}

% ===========================================================
% 19. ECOSYSTEM
% ===========================================================
\section{SPM-Kit ecosystem relationship}

\begin{itemize}
  \item \textbf{SPM-Kit Core} reads demonstrated variants and performs numerical analysis.
  \item \textbf{Fathom} invokes Core through an interactive workspace.
  \item \textbf{Data Hunter} records public candidates and provenance; discovery is not validation.
  \item \textbf{Phantoms} generates deterministic truth-bearing synthetic fixtures independently of the analyzer.
  \item \textbf{Validation} runs frozen black-box campaigns against installed public interfaces.
\end{itemize}

Artifacts move through declared files, manifests, hashes, and commands. Not every located
candidate is used, accepted, redistributable, or scientifically suitable, and there is no
automatic pipeline from every Data Hunter record into Validation.

% ===========================================================
% 20. CITATION AND LICENSE
% ===========================================================
\section{Citation, license, and contributions}

\subsection{How to cite}

\begin{lstlisting}[language=bash]
@software{spmkit2026,
  author       = {Jos\'e Labarca Baeza},
  title        = {spmkit: Open-source AFM/KPFM analyser
                  for scanning probe microscopy},
  year         = {2026},
  version      = {0.1.4},
  doi          = {10.5281/zenodo.21303280},
  url          = {https://github.com/kegouro/spmkit},
}
\end{lstlisting}

\subsection{License}
MIT. See \texttt{LICENSE} in the repository root.

\subsection{Contributing}
Issues and pull requests at \url{https://github.com/kegouro/spmkit/issues}.
See \texttt{CONTRIBUTING.md}.

\subsection{Funding}
This open-source project has been developed without dedicated institutional funding.

\subsection{Project status}
\textbf{Alpha (source \texttt{0.1.5.dev0}; GitHub release \texttt{0.1.4}; PyPI
\texttt{0.1.2}).} APIs may change before 1.0. Evidence and compatibility remain scoped.

% ===========================================================
% 21. ACKNOWLEDGEMENTS
% ===========================================================
\section{Acknowledgements}

Jos\'e Labarca Baeza is the creator, author, and lead developer of SPM-Kit.

\subsection{English}

Tomás Corrales and the SPM Lab at Universidad Técnica Federico Santa María provided selected experimental datasets and laboratory context during the development and evaluation of SPM-Kit.

María Saavedra Fredes and Benjamin Schleyer helped locate and share candidate datasets for the validation campaigns.

\subsection{Español}

Tomás Corrales y el SPM Lab de la Universidad Técnica Federico Santa María proporcionaron datasets experimentales seleccionados y contexto de laboratorio durante el desarrollo y la evaluación de SPM-Kit.

María Saavedra Fredes y Benjamin Schleyer ayudaron a localizar y compartir datasets candidatos para las campañas de validación.

These acknowledgements do not imply that every located dataset was used, accepted,
redistributable, or scientifically suitable. They do not constitute software authorship,
institutional ownership, or endorsement.

% ===========================================================
% APPENDICES
% ===========================================================
\appendix

\section{Quick reference card}

\begin{lstlisting}[language=bash]
# Install the documented source
python -m pip install "spmkit[gui] @ git+https://github.com/kegouro/spmkit@main"

# Inspect
spmkit info scan.nid

# Analyse
spmkit roughness scan.nid -c Z-Axis --level plane
spmkit psd scan.nid -c Z-Axis
spmkit analyze scan.nid -o ./results --tip-wf 4.8

# Figures
spmkit figure scan.nid -o topo.png --colormap batlow

# Convert
spmkit convert scan.nid scan.gwy

# Verify
spmkit verify scan.nid

# Force
spmkit forcecurve curve.jpk-force --model dmt --tip-radius 2e-8
spmkit forcemap volume.nid --fast -f maps.png
spmkit forcereport volume.nid -o report
spmkit forceexport volume.nid -o ./export
spmkit fbatch ./curves/ -o batch.csv

# Evaporation
spmkit evaporation ./tuning/ -k 0.3 -o evap.csv

# Grains
spmkit grains scan.nid --min-size 10

# GUI
spmkit gui [file]
\end{lstlisting}

\section{CLI command index}

\begin{longtable}{@{}>{\ttfamily}l l l@{}}
\toprule
Command & Category & Input \\
\midrule
info & Inspection & .nid, .nhf \\
roughness & Image & .nid, .nhf, .gwy \\
psd & Spectral & .nid, .nhf, .gwy \\
analyze & Pipeline & .nid, .nhf, .gwy \\
nanomech & Force & .nid (spectroscopy) \\
grains & Image & .nid, .nhf, .gwy \\
batch & Batch & Folder of SPM \\
evaporation & Resonance & Folder of .nid tuning \\
figure & Figure & .nid, .nhf, .gwy \\
convert & Format & .nid \\
verify & Audit & .nid \\
gui & GUI & --- (optional file) \\
workspace & GUI & Force-curve file \\
forcecurve & Force & .jpk-force, .nid \\
forcemap & Force & .nid (force-volume) \\
forcereport & Force & .nid (force-volume) \\
forceexport & Force & .nid (force-volume) \\
fbatch & Batch & Folder of force curves \\
jkr & Force (exp.) & Calibrated curve \\
\bottomrule
\end{longtable}

\section{Perspective index}

\begin{longtable}{@{}>{\ttfamily}l l l l@{}}
\toprule
Key & Label & Module & Type \\
\midrule
image & Imagen & image & Image \\
grains & Granos & image & Image \\
spectral & Espectral & image & Image \\
resonance & Sintonía térmica & image & Image \\
evaporation & Evaporación & image & Image \\
force & Curva de fuerza & force & Force \\
smfs & SMFS & force & Force \\
map & Mapa & force & Force \\
batch & Batch & force & Force \\
figure & Figura & figure & Figure \\
view3d & Vista 3D & view3d & Image \\
simulator & Simulador & simulator & Educational \\
\bottomrule
\end{longtable}

\section{Supported formats matrix}

\begin{longtable}{@{}l l l l l@{}}
\toprule
Extension & Source & Type & Access & Status \\
\midrule
.nid & NanoSurf & Image + force & Read & Implemented, scoped \\
.nhf & NanoSurf HDF5 & Images & Read & Experimental \\
.gwy & Gwyddion & Images & Read+Write & Interoperability \\
.spm/.00N & Bruker/Nanoscope & Images & Read & \textbf{Partial LEVEL 2} \\
.jpk-force/.jpk & JPK & Single curve & Read & Implemented, scoped \\
JPK tagged TIFF & JPK & Force data & Read & Experimental \\
.ibw/.h5/.tab & Adapter sources & Force data & Read & Experimental \\
.npz & Phantoms & Images & Read & Declared bundle \\
\bottomrule
\end{longtable}

\section{Keyboard shortcuts}

\begin{tabularx}{\linewidth}{@{}>{\ttfamily}l X@{}}
\toprule
Shortcut & Action \\
\midrule
Ctrl+K & Command palette \\
Ctrl+O & Open file \\
Ctrl+S & Save project \\
Ctrl+M & Calculate map \\
Ctrl+Shift+R & Generate report \\
Ctrl+Shift+L & Toggle theme \\
Ctrl+Shift+A & Appearance settings \\
Ctrl+Left / Ctrl+Right & Previous / next curve \\
Ctrl+1 ... Ctrl+5 & Tab switch (legacy app) \\
\bottomrule
\end{tabularx}

\vfill

\begin{center}
{\small --- End of SPM-Kit {\textperiodcentered} Fathom Usage Guide --- source 0.1.5.dev0 {\textperiodcentered} GitHub 0.1.4 {\textperiodcentered} PyPI 0.1.2 ---}
\end{center}

\end{document}
